¿Por qué Kali Linux es tan popular en seguridad informática?

Kali Linux es una distribución de Linux popular en el área de seguridad informática debido a la naturaleza con la que fue diseñada, pues se construyó con la idea de ser una distribución para realizar pentesting y auditorías de seguridad, por lo que contiene por defecto múltiples herramientas que facilitan esta tarea. Además, por defecto Kali Linux incluye hooks de systemd que deshabilitan los servicios de red por defecto, lo que permite probar e instalar servicios a la vez que se garantiza la seguridad de la distribución.

Kali Linux utiliza un Kernel de Linux modificado para la inyección inalámbrica y un conjunto mínimo y fiable de repositorios para mantener la seguridad del sistema~\cite{kali_docs}.

Fuera de las bondades técnicas que tiene Kali Linux, esta distribución es popular en el área de seguridad informática por la gran cantidad de herramientas de seguridad diseñadas para sistemas Linux, de las que muchas se incorporan por defecto, además de ser una distribución relativamente sencilla de entender para usuarios con una mínima experiencia en sistemas Linux y ser un proyecto bien mantenido y documentado.
